% !PRINTED 285
\weberpage{285}

%<<CONT
then by \S.\,83, V. the number of sign changes in
\weq{2}{1,\ \varDelta_1(\alpha), \varDelta_2(\alpha) \dots \varDelta_n(\alpha)}
is equal to the number $N_\alpha$ of negative squares in $H$, and if we therefore count the corresponding number $N_\beta$ in
\[
1,\ \varDelta_1(\beta), \varDelta_2(\beta) \dots \varDelta_n(\beta)
\]
then $N_\beta - N_\alpha$ is the number of roots of $f(x)$ lying between $\alpha$ and $\beta$.

For the cubic equation the expressions follow from the end of the preceding section. We form the chain
\[
a_0,\quad \frac{1}{a_0} H_{2,2},\quad \frac{1}{a_0}\left(H_{2,2} H_{1,1} - H_{1,2}^2\right),\quad \frac{1}{a_0} \varDelta,
\]
which evidently has the same number of sign changes as (2), and thus obtain the four functions
\[
\begin{aligned}
&a_0,\qquad -3 a_0 \alpha - a_1,\\
&2\alpha^2 a_0(a_1^2 - 3a_0 a_2) + \alpha(2a_1^3 + 9a_0^2a_3 - 7a_0a_1a_2)\\
&\quad + 3a_0a_1a_3 + a_1^2a_2 - 4a_0a_2^2,\qquad -f(\alpha)\,D.
\end{aligned}
\]
If as a test we take $\alpha = -\infty$, $\beta = +\infty$, then, when moreover $a_0 = 1$ is set, we obtain the determination of signs
\[
\begin{aligned}
1, &\quad +1, \quad (a_1^2 - 3a_2), \quad D,\\
1, &\quad -1, \quad (a_1^2 - 3a_2), \quad -D.
\end{aligned}
\]
In the enumeration one must observe that when $D$ is positive, $a_1^2 - 3a_2$ cannot be negative [\S.\,47, (8)].

\begin{center}
\S.\,93.

\gesperrt{Outlines of the theory of characteristics}\footnote{Monatsberichte der Berliner Akademie, March and August 1869, February 1873, February 1878.}).
\end{center}

The theorems of \gesperrt{Sturm} are generalised in an extraordinary manner by the theory of characteristics of \gesperrt{Kronecker}, which pursues the same aim as \gesperrt{Sturm}'s theorem for systems of equations in arbitrarily many variables. We confine ourselves here to the consideration of the simplest case, which we intend to apply to the enclosure of the complex roots of an equation; moreover, we here avail ourselves un-
%CONT>>
